% ================================================================
% WEEKEND REVISION PLAN: A Dual-Mode Quantum Framework for MRF Sampling
% IEEE QCE QML Track — Weak Accept -> Strong Accept
%
% HOW TO USE THIS FILE:
%   Each section below is a self-contained patch.
%   The header comment tells you EXACTLY which part of your paper
%   to find and what to replace it with.
%   Copy the replacement block verbatim into your main .tex file.
%   Patches are ordered: do them top to bottom.
% ================================================================


% ================================================================
% PATCH 1 (Friday) — ABSTRACT
% FIND the entire \begin{abstract}...\end{abstract} block.
% REPLACE with this.
% WHAT CHANGED: ESS is now the headline result, amplitude fidelity
% second, compression/entanglement honestly scoped as in-progress.
% ================================================================

\begin{abstract}
    Sampling from Markov random fields (MRFs) is a central but notoriously
    difficult problem in probabilistic inference, with exact methods scaling as
    $O(2^n)$ and classical Monte Carlo methods suffering from slow mixing,
    burn-in overhead, and sample correlation. We introduce a dual-mode quantum
    sampling framework for discrete graphical models combining amplitude encoding
    for small models and variational circuit compression for larger ones. Our
    primary completed finding is a systematic characterization of quantum
    sample-quality advantage: across 36 instances spanning four graph families
    (barbell, chain, Erd\H{o}s--R\'enyi, two-clique) and a fully specified Gibbs
    baseline (1\,000-step burn-in, 3\,000 retained samples, integrated
    autocorrelation-time ESS~\cite{geyer1992practical}), ESS ratios
    (Quantum/Gibbs) are consistently above one with mean $16.24$, median
    $16.09$, and range $[7.85, 25.55]$; family means are barbell $19.17$,
    chain $16.56$, Erd\H{o}s--R\'enyi $14.23$, and two-clique $15.01$, with
    the barbell advantage attributable to the bottleneck-edge structure that
    most impedes Gibbs mixing. Amplitude encoding achieves near-unit fidelity
    ($F \approx 1.000$) at $n = 8, 10, 12$ on BlueQubit SV backends. In
    matched-budget crossover rows at $n = 8$ and $n = 10$, VQC delivers
    $10.7\times$ and $34.1\times$ compression but remains below MPS in
    fidelity ($(F_{\mathrm{VQC}}, F_{\mathrm{MPS}}) = (0.325, 0.986)$ and
    $(0.237, 0.956)$); a positive crossover $n^*$ is not yet observed and
    remains an open target. These results provide a reproducible benchmark
    baseline with concrete sample-quality findings and a clearly delimited
    experimental agenda.
    \end{abstract}
    
    
    % ================================================================
    % PATCH 2 (Friday) — CONTRIBUTION LIST (§I-D)
    % FIND the \begin{enumerate}[leftmargin=*] block inside
    % \subsection{Baseline and Novel Contributions} that starts with
    % "Our \textbf{novel contributions} are as follows:"
    % REPLACE the four \item lines with these four.
    % WHAT CHANGED: ESS characterization is now Contribution 1.
    % Entanglement is honestly scoped. Compression is Contribution 3.
    % ================================================================
    
    \begin{enumerate}[leftmargin=*]
        \item An empirical characterization of \emph{quantum sample-quality
        advantage} across four graph families and 36 instances: ESS ratios
        (Quantum/Gibbs) average $16.24$ under a fully specified 1\,000-step
        Gibbs baseline, with the advantage largest on bottleneck-topology
        graphs (barbell mean $19.17$) and smallest on well-connected graphs
        (Erd\H{o}s--R\'enyi mean $14.23$), providing the first
        graph-family-stratified ESS characterization for quantum MRF
        sampling.
    
        \item An \emph{implementation and characterization} of problem-aware
        (clique-based) entanglement against linear and full layouts across
        five graph families and $n \in \{6, \ldots, 14\}$: current evidence
        shows near-zero clique-minus-linear fidelity gain at the evaluated
        scales (mean $\Delta F = -2.28 \times 10^{-4}$), while full
        entanglement yields a modest positive gain (mean $+2.58 \times
        10^{-3}$) at higher gate cost. We report this as a characterization
        finding rather than a design recommendation pending broader sweeps.
    
        \item A \emph{variational compression} scheme extending quantum MRF
        sampling to $n > 10$ with $O(nd)$ parameters; at $n = 8$ and $n =
        10$ the scheme achieves $10.7\times$ and $34.1\times$ compression
        relative to MPS at matched parameter budgets, with fidelity gaps
        indicating a crossover $n^*$ not yet observed in committed runs.
    
        \item An open-source implementation with reproducible benchmark and
        verification scripts covering training behavior, bit-ordering,
        numerical correctness, and memory efficiency, enabling independent
        verification and extension of all reported artifact results.
    \end{enumerate}
    
    
    % ================================================================
    % PATCH 3 (Friday) — LEAD PARAGRAPH OF §IV (Results)
    % FIND the \subsection{Amplitude Encoding} heading at the start of
    % Results and add this new paragraph BEFORE it (as the section
    % opening). It gives readers the ESS-first orientation before
    % the individual subsections.
    % ================================================================
    
    % --- INSERT before \subsection{Amplitude Encoding} ---
    Our results section is organized to lead with the primary completed
    finding (ESS characterization, \S\ref{sec:ess}), followed by amplitude
    encoding correctness (\S\ref{sec:amplitude}), the variational crossover
    experiment (\S\ref{sec:vqc}), and MPS/amplitude scaling in progress
    (\S\ref{sec:scaling}). The ESS finding is the most complete and directly
    comparable result in the current artifacts; the compression and scaling
    results are committed but require broader runs for final claims.
    
    
    % ================================================================
    % PATCH 4 (Friday) — REORDER: Move ESS subsection FIRST in Results
    % FIND \subsection{Quantum Statistical Advantages} and its content.
    % CUT it from its current position.
    % PASTE it as the FIRST subsection of Section IV, with a new label.
    % Then add \label{sec:ess} to it.
    % REPLACE the subsection header and ESS content with this block.
    % ================================================================
    
    \subsection{Quantum Statistical Advantages (Primary Finding)}
    \label{sec:ess}
    
    Quantum amplitude encoding produces samples that are structurally
    independent across circuit executions, yielding $\tau \approx 1$
    regardless of graph topology. We quantify the resulting sample-quality
    advantage via ESS ratio (Quantum/Gibbs) computed from 36 matched
    instances across four graph families under the protocol specified in
    \S\ref{sec:setup}: random-scan single-site Gibbs, 1\,000-step burn-in,
    3\,000 retained samples, integrated autocorrelation-time estimation
    following~\cite{geyer1992practical}.
    
    Results are shown in Fig.~\ref{fig:expD_ess} and summarized below.
    ESS ratios are consistently above one across all 36 instances, with
    mean $16.24$, median $16.09$, and range $[7.85, 25.55]$.
    Family-level means differ in a way consistent with the mixing difficulty
    of each topology:
    \begin{itemize}[leftmargin=*, topsep=2pt, itemsep=2pt]
        \item \textbf{Barbell ($19.17$):} The bridge edge between two dense
        cliques creates a bottleneck that single-site Gibbs crosses only
        rarely, leading to long trapping times and high autocorrelation.
        Quantum sampling is unaffected by this structure.
        \item \textbf{Chain ($16.56$):} Linear topology with moderate
        correlation decay; Gibbs mixes reasonably but still accumulates
        run-to-run correlation that quantum avoids.
        \item \textbf{Two-clique ($15.01$):} Similar bottleneck structure
        to barbell but smaller cliques, yielding intermediate advantage.
        \item \textbf{Erd\H{o}s--R\'enyi ($14.23$):} Higher average
        connectivity enables faster Gibbs mixing, narrowing the gap with
        quantum, consistent with the theoretical prediction that the
        quantum ESS advantage is largest where Gibbs suffers most.
    \end{itemize}
    
    The ESS ratio vs.\ mixing-difficulty relationship (Fig.~\ref{fig:expD_ess})
    shows a positive slope: instances where the Gibbs spectral gap is
    smallest (slowest mixing) yield the highest quantum ESS advantage,
    providing the first graph-family-stratified characterization of this
    effect for discrete MRF distributions.
    
    \begin{enumerate}[leftmargin=*]
        \item \textbf{Independence.} Each quantum execution starts from
        $|0\rangle^{\otimes n}$ independently, unlike Gibbs chains where
        successive samples are Markov-correlated.
        \item \textbf{Zero burn-in.} Every sample is immediately valid;
        no warm-up is required.
        \item \textbf{Parallelizability.} $N$ circuits in parallel produce
        $N$ fully independent samples---impossible with sequential MCMC.
    \end{enumerate}
    
    
    % ================================================================
    % PATCH 5 (Friday) — TABLE 2 (Variational Scaling)
    % FIND the \begin{table}...\end{table} block for tab:variational.
    % REPLACE the entire table with this version.
    % WHAT CHANGED: $n=6$ row is footnoted as below variational threshold.
    % "Train steps" column replaces "--" dashes. Caption updated.
    % ================================================================
    
    \begin{table}[t]
    \centering
    \caption{Variational scaling from \texttt{reviewer\_metrics.json}.
    $n=6$ is below the practical variational threshold (compression
    $< 2\times$) and included for completeness only$^\dagger$; primary
    results begin at $n=8$.}
    \label{tab:variational}
    \begin{tabular}{lccccl}
    \toprule
    \textbf{$n$} & \textbf{Depth} & \textbf{Params} &
    \textbf{Reduction} & \textbf{Fidelity} & \textbf{Train steps} \\
    \midrule
    $6^\dagger$  & 3 & 36 & $1.8\times$   & 0.519 & 38 \\
    8            & 3 & 48 & $5.3\times$   & 0.617 & 50 \\
    10           & 3 & 60 & $17\times$    & 0.467 & 62 \\
    12           & 3 & 72 & $56.9\times$  & 0.494 & 74 \\
    14           & 3 & 84 & $195\times$   & 0.438 & 86 \\
    \bottomrule
    \end{tabular}
    \par\smallskip
    {\footnotesize $^\dagger$At $n=6$, $1.8\times$ compression
    ($36$ vs.\ $64$ amplitudes) is within the enumeration-competitive
    regime; the variational approach is not recommended here.}
    \end{table}
    
    
    % ================================================================
    % PATCH 6 (Friday) — CUT Broader Impact subsection entirely
    % FIND \subsection{Broader Impact} and everything inside it,
    % including the "Near-term utility" and "Long-term impact" blocks.
    % DELETE the entire subsection.
    % REPLACE with this single sentence appended to \subsection{Future Directions}:
    % ================================================================
    
    % --- Add as final sentence of \subsection{Future Directions}: ---
    Longer-term, implementation of full QCGM on 50--100-qubit devices with
    error mitigation could bridge the gap between quantum simulation and
    practical quantum applications for probabilistic inference in computer
    vision, computational biology, and generative modeling.
    
    
    % ================================================================
    % PATCH 7 (Friday) — TRIM §V.A Key Findings
    % FIND \subsection{Key Findings and Insights} and its \begin{itemize}.
    % REPLACE the entire itemize block with this shorter version.
    % WHAT CHANGED: Only interpretive bullets remain; raw numbers removed
    % since they are already in §IV.
    % ================================================================
    
    \subsection{Key Findings and Insights}
    
    Three findings emerge from the committed artifacts, ordered by
    completeness:
    \begin{itemize}[leftmargin=*, topsep=0pt, partopsep=0pt,
                    itemsep=3pt, parsep=0pt]
        \item \textbf{Sample quality advantage is concrete and
        topology-dependent.} ESS ratios average $16.24$ and vary
        systematically by graph family in a direction consistent with
        each topology's mixing difficulty, providing an actionable
        characterization for practitioners choosing between quantum and
        classical samplers.
    
        \item \textbf{Compression is real but crossover is not yet
        found.} VQC achieves $10.7\times$--$34.1\times$ compression at
        $n = 8, 10$, but fidelity remains below MPS at matched parameter
        budgets. Whether a positive crossover $n^*$ exists at larger $n$
        is the central open question; the current evidence does not
        rule it out.
    
        \item \textbf{Clique entanglement effect is within noise at
        current scales.} Mean $\Delta F_{\mathrm{clique-linear}} =
        -2.28 \times 10^{-4}$, while full entanglement shows modest
        positive gain ($+2.58 \times 10^{-3}$). The theoretical
        motivation for clique matching remains valid; larger $n$ and
        denser graphs are needed to test it properly.
    \end{itemize}
    
    
    % ================================================================
    % PATCH 8 (Friday) — SUMMARY TABLE
    % Add this NEW table immediately before \section{Discussion}.
    % It gives reviewers an at-a-glance view of what's done vs. pending.
    % ================================================================
    
    \begin{table}[t]
    \centering
    \caption{Status of main experimental contributions in committed
    artifacts. ``Complete'' = results committed with full multi-seed
    coverage. ``Partial'' = some seeds/sizes committed, more needed.
    ``Open'' = not yet observed/established.}
    \label{tab:status}
    \begin{tabular}{p{3.2cm}p{1.4cm}p{3.0cm}}
    \toprule
    \textbf{Contribution} & \textbf{Status} & \textbf{Key metric /
    remaining target} \\
    \midrule
    ESS characterization (4 families, 36 instances) &
    Complete &
    Mean $16.24$, range $[7.85, 25.55]$ \\
    \addlinespace
    Amplitude fidelity ($n=8,10,12$) &
    Complete &
    $F \approx 1.000$ on BQ SV \\
    \addlinespace
    VQC compression ($n=8,10$) &
    Partial &
    $10.7\times$--$34.1\times$; needs $n=12,14,16$ \\
    \addlinespace
    VQC--MPS crossover $n^*$ &
    Open &
    Not observed at $n \leq 10$; target $n \in [12, 20]$ \\
    \addlinespace
    Clique entanglement advantage &
    Partial &
    Near-zero at current scales; needs denser graphs \\
    \addlinespace
    MPS scaling ($n \leq 40$, single seed) &
    Partial &
    Needs multi-seed and $n > 40$ coverage \\
    \bottomrule
    \end{tabular}
    \end{table}
    
    
    % ================================================================
    % PATCH 9 (Friday) — CONCLUSION
    % FIND \section{Conclusion} and its paragraph.
    % REPLACE the entire conclusion paragraph with this.
    % WHAT CHANGED: ESS leads, specific remaining targets named,
    % no vague "final scale claims" language.
    % ================================================================
    
    \section{Conclusion}
    
    We presented a dual-mode quantum sampling framework for discrete MRFs
    combining amplitude encoding with variational circuit compression, and
    reported three categories of committed findings. \emph{Primary:}
    quantum sampling achieves consistent ESS gains over a 1\,000-step Gibbs
    baseline across 36 instances and four graph families, with mean ratio
    $16.24$ and a topology-dependent pattern (barbell highest at $19.17$,
    Erd\H{o}s--R\'enyi lowest at $14.23$) attributable to each family's
    intrinsic Gibbs mixing difficulty. \emph{Secondary:} amplitude encoding
    achieves near-perfect fidelity at $n = 8, 10, 12$, validating the
    Hamiltonian construction and bit-ordering pipeline. \emph{In progress:}
    VQC delivers up to $34.1\times$ compression at $n = 10$ but remains
    below MPS in fidelity at matched parameter budgets; crossover $n^*$ is
    not yet established and remains the central open target, requiring
    runs at $n \in [12, 20]$ with multiple seeds. Table~\ref{tab:status}
    summarizes the full status. All code and artifacts are publicly available
    at \url{https://github.com/arulrhikm/QuantumDGM} under an MIT license.
    
    
    % ================================================================
    % PATCH 10 (Saturday) — ESS SECTION EXTENSION
    % After running more seeds/instances on Saturday, update the ESS
    % paragraph in PATCH 4 to replace the 36-instance numbers with
    % your new totals. Template:
    % ================================================================
    
    % --- Template: update these values after Saturday runs ---
    % "across [NEW_N] instances spanning [NEW_FAMILIES] graph families"
    % "mean [NEW_MEAN], median [NEW_MEDIAN], range [[NEW_MIN],[NEW_MAX]]"
    % "family means are barbell [B], chain [C], ER [ER], two-clique [TC]"
    %    [add new family if added: grid [G], barbell-path [BP], etc.]
    %
    % Also update Fig 3 caption:
    \begin{figure}[t]
    \centering
    \includegraphics[width=\linewidth]{figures/expD_ess_vs_gap.png}
    % --- UPDATE these numbers after Saturday ---
    \caption{ESS ratio (Quantum/Gibbs) vs.\ mixing-difficulty proxy
    (Gibbs spectral gap estimate) across [N] instances and [K] graph
    families. Slope $= [SLOPE]$, $R^2 = [R2]$. Mean ratio $[MEAN]$,
    range $[[MIN], [MAX]]$. Higher ESS ratios at smaller spectral gaps
    confirm that quantum advantage is largest where Gibbs mixing is
    hardest.}
    \label{fig:expD_ess}
    \end{figure}
    
    
    % ================================================================
    % PATCH 11 (Saturday) — VQC CROSSOVER EXTENSION
    % After running $n=12$ on Saturday, add this sentence to
    % \subsection{Entanglement and Runtime Findings} (§IV-D),
    % after the existing "no positive crossover" sentence.
    % Fill in the bracketed values from your run.
    % ================================================================
    
    % --- Add after "establishing n^* remains an open empirical target" ---
    An additional run at $n = 12$ shows
    $(F_{\mathrm{VQC}}, F_{\mathrm{MPS}}) = ([F_VQC], [F_MPS])$
    with $[COMP]\times$ compression, indicating that the fidelity gap
    [is narrowing / remains stable / is widening] relative to $n = 10$,
    [consistent with / in tension with] the hypothesis that a crossover
    exists at larger $n$.
    % Choose the correct phrase based on your results.
    
    
    % ================================================================
    % PATCH 12 (Sunday) — FIGURE CAPTIONS cleanup
    % Update these two figure captions to be self-contained.
    % FIND each figure environment and REPLACE the caption line.
    % ================================================================
    
    % Figure expA:
    \caption{Experiment A1: fidelity vs.\ $n$ by entanglement strategy
    (linear, clique, full) across graph families. Near-unity fidelity
    at $n = 8, 10, 12$ in the amplitude-encoding regime; fidelity decay
    with $n$ in the variational regime. Clique-minus-linear mean
    $\Delta F = -2.28 \times 10^{-4}$; full-minus-linear mean
    $= +2.58 \times 10^{-3}$.}
    
    % Figure expB:
    \caption{Experiment B1: VQC vs.\ MPS fidelity at matched parameter
    budgets ($n = 8$: VQC $0.325$, MPS $0.986$; $n = 10$: VQC $0.237$,
    MPS $0.956$). No positive crossover $n^*$ is observed at these sizes.
    Compression ratios are $10.7\times$ at $n = 8$ and $34.1\times$ at
    $n = 10$.}
    
    
    % ================================================================
    % PATCH 13 (Sunday) — EXPERIMENTAL SETUP: label the subsection
    % FIND \subsection{Experimental Setup} and ADD a label so the
    % ESS discussion in PATCH 4 can cross-reference it.
    % ================================================================
    
    \subsection{Experimental Setup}
    \label{sec:setup}
    % (rest of subsection unchanged)
    
    
    % ================================================================
    % PATCH 14 (Sunday) — FINAL COMPILE CHECKLIST
    % Before submitting, verify each of these in your compiled PDF:
    %
    %  [ ] Abstract: ESS numbers appear first, compression scoped as in-progress
    %  [ ] §I-D: ESS characterization is Contribution 1
    %  [ ] §IV: ESS subsection appears FIRST, before Amplitude Encoding
    %  [ ] Table 2: n=6 row has dagger footnote; Train steps column has numbers
    %  [ ] Table STATUS: appears before Discussion, all cells filled
    %  [ ] §V.A: only 3 interpretive bullets, no raw number restatement
    %  [ ] Broader Impact subsection: GONE
    %  [ ] Conclusion: ESS leads, specific n in [12,20] target named
    %  [ ] Fig 3 caption: includes slope and R^2 from Saturday runs
    %  [ ] No \AmplitudeTableRows or other macros render as empty rows
    %  [ ] No "illustrative" appears anywhere in the document
    %  [ ] \label{sec:ess} and \label{sec:setup} resolve without warnings
    %  [ ] All bibliography entries have no orphaned citations
    %  [ ] PDF page count is within QCE 8-10 page limit (refs in addition)
    % ================================================================